%% kapa: exam=1 round=37 year=2026 period=1 form=A subject=real_estate subject_ko=부동산학원론 seq=24 number=104
%% source: Q-Net 감정평가사 기출문제 (cst00309) artlSeq=5258924 file=raw/1008/5258924/2026년 제37회 감정평가사 1차 1교시.pdf
%% conversion: kordoc 4.24.51 -> tools/gen1.py (auto); review=proofread-sonnet; flags=none
\begin{question}{104}
부동산 투자분석기법에 관한 설명으로 옳지 않은 것은? (단, 주어진 조건에 한함)
\begin{choices}
  \item 수익성지수(PI)가 ‘1’보다 큰 경우 투자가치가 있다고 판단한다.
  \item 내부수익률(IRR)이 요구수익률보다 큰 경우 투자가치가 있다고 판단한다.
  \item 순현재가치(NPV)법은 자본비용으로 재투자한다고 가정하며 가치가산의 원리가 적용된다.
  \item 다른 조건이 일정하다면 승수법에서는 승수가 클수록 더 좋은 투자안이다.
  \item 투자규모에 차이가 나는 상호배타적인 투자안을 검토하는 경우 순현재가치(NPV)법과 내부수익률(IRR)법을 통한 의사결정이 달라질 수 있다.
\end{choices}
\end{question}
\answer{4}
